% BEGIN REV09_MEMORIA_PROFUNDA
\subsection{The deep coordinate and integer recovery}
\label{corpus9:xii:memoria-profunda}

The six-coordinate chart retains the expressed residue and the integer
quotient that determines its prolongation. We fix the integer
representative of \(L_0\) displayed in the regional construction:
\begin{equation}
 A_H=\begin{pmatrix}
 2&2&2&1&2&1\\
 2&1&2&2&1&1\\
 1&0&0&1&0&2\\
 0&0&1&1&1&0\\
 2&2&2&0&2&1\\
 2&1&2&1&0&0
 \end{pmatrix},\qquad \det A_H=1.
 \label{corpus9:xii:matriz-profunda}
\end{equation}
Its reduction modulo three equals \(L_0\). The choice of integer
representative is part of the chart: in addition to that reduction, it
fixes the action on the deep quotients. For row vectors, define
\begin{equation}
 \begin{aligned}
 \mathsf{Div}_H:\mathbb Z^6&\longrightarrow
       \{0,1,2\}^6\times\mathbb Z^6,\\
 x_H&\longmapsto(u_H,x_H^+),\\
 y_H&=x_HA_H,\qquad u_H=[y_H]_3,\qquad
 x_H^+=(y_H-u_H)/3.
 \end{aligned}
 \label{corpus9:xii:division-profunda}
\end{equation}
Reduction \([\,\cdot\,]_3\) is performed componentwise with
representatives \(0,1,2\), including for negative integer components.

\begin{proposition}[Exact inversion of residue and quotient]
\label{corpus9:xii:inversion-profunda}
The map \(\mathsf{Div}_H\) is bijective. Its inverse is
\begin{equation}
 x_H=(u_H+3x_H^+)A_H^{-1}.
 \label{corpus9:xii:inversa-profunda}
\end{equation}
For a trajectory in this chart,
\(x_jA_H=u_j+3x_{j+1}\), reconstruction at depth
\(T\geq1\) satisfies
\begin{equation}
 x_0=\sum_{j=0}^{T-1}3^ju_jA_H^{-(j+1)}
                 +3^Tx_TA_H^{-T}.
 \label{corpus9:xii:reconstruccion-profunda}
\end{equation}
\end{proposition}
\begin{proof}
The determinant in \eqref{corpus9:xii:matriz-profunda} is one, so the
adjugate formula gives \(A_H^{-1}\in M_6(\mathbb Z)\).
Euclidean division of each component of \(x_HA_H\) determines a
unique pair \((u_H,x_H^+)\). Conversely, for any
\((u,q)\in\{0,1,2\}^6\times\mathbb Z^6\), the vector
\(x=(u+3q)A_H^{-1}\) is integral and satisfies \(xA_H=u+3q\).
Its residue is exactly \(u\), and its quotient is \(q\). This proves
both inverse compositions. Successive substitution of
\(x_j=(u_j+3x_{j+1})A_H^{-1}\), for \(j=0,\ldots,T-1\), gives
\eqref{corpus9:xii:reconstruccion-profunda}; the final term retains
the terminal coordinate remaining after the observed prefix.
\end{proof}

The ternary orientation is read subsequently through
\(0\mapsto0\), \(1\mapsto+1\), and \(2\mapsto-1\), retaining
the residue \(u_H\) and the quotient in
\eqref{corpus9:xii:division-profunda}. The other state fields preserve
their own records of orientation, carry, route, and boundary.
Proposition~\ref{prop:ch-hensel-bruto} establishes the arithmetic tower
for \(p=3\), \(d=6\), and \(A=A_H\), together with proofs of
compatibility and passage to the limit. Realization of the generated
prolongations also uses the APP--TRIT--TPK compatibility conditions
specified in Proposition~\ref{prop:cobertura-729-seccion}.

The residual block also retains its index
\[
 \nu(u_H)=\sum_{j=1}^{6}(u_H)_j3^{6-j}\in\{0,\ldots,728\}.
\]
This index enters the prefixes through
\(N'=729N+\nu(u_H)\). The mixed chart in
Section~\ref{act21:subsec:df-cilindros-mixtos} simultaneously retains
the ternary and decimal indices and the integer boundary quantities
\(A(s),B(s)\). Their recurrences and compatibility with truncation are
proved in Proposition~\ref{act21:prop:df-rec-frontera} and
Theorem~\ref{act21:thm:df-refinamiento-generado}. Integer-coordinate
recovery is thereby linked to the exact comparison of the cylinders
that express its prefixes.
% END REV09_MEMORIA_PROFUNDA
